\section{问题一完整默认参数附表}
\begin{table}[htbp]
\centering
\caption{问题一完整默认参数附表}
\label{tab:q1_param_full}
\scriptsize
\setlength{\tabcolsep}{2.5pt}
\begin{tabularx}{0.93\textwidth}{>{\raggedright\arraybackslash}p{2.0cm} >{\raggedright\arraybackslash}p{1.6cm} >{\raggedright\arraybackslash}X >{\raggedright\arraybackslash}p{1.6cm}}
\toprule
参数 & 默认值 & 依据/说明 & 角色 \\
\midrule
\multicolumn{4}{l}{\textbf{资源增长与归一化基准}} \\
$r_1,r_2,r_3,r_4$ & \makecell[l]{$1.15,0.95,$\\$0.82,0.72$} & 归一化相对恢复速度排序，表示水草快于浮游植物、浮游植物快于浮游动物、浮游动物快于底栖饵料；用于维持营养级时间尺度分离。 & 时间尺度约束 \\
$K_1,K_2,K_3,K_4$ & $(1,1,1,1)$ & 统一固定为归一化承载上界，不对应河段实测承载量。 & 归一化基准 \\
\midrule
\multicolumn{4}{l}{\textbf{摄食耦合与饱和响应}} \\
$\alpha_1,\alpha_2,\alpha_3,\alpha_4$ & \makecell[l]{$1.22,0.92,$\\$0.84,0.76$} & Holling II 型摄食耦合强度，按 $j=1,2,3,4$ 分别对应水草--草鱼、浮游植物--鲢鱼、浮游动物--鳙鱼、底栖饵料--青鱼\cite{holling1959}。 & 结构耦合 \\
$h_1,h_2,h_3,h_4$ & \makecell[l]{$0.55,0.50,$\\$0.48,0.45$} & 饱和摄食拐点参数，使系统在归一化尺度下保持有界响应。 & 饱和约束 \\
\midrule
\multicolumn{4}{l}{\textbf{转化效率与自然损失}} \\
$e_1,e_2,e_3,e_4$ & \makecell[l]{$0.43,0.35,$\\$0.33,0.30$} & 归一化转化效率；沿“草鱼--水草、鲢鱼--浮游植物、鳙鱼--浮游动物、青鱼--底栖饵料”顺序设置为递减，避免消费者收益超过资源层可供给量级。 & 能量转化 \\
$m_1,m_2,m_3,m_4$ & \makecell[l]{$0.16,0.14,$\\$0.13,0.12$} & 归一化自然损失系数，与转化效率共同维持消费者层在禁渔和非禁渔情景下的数量级稳定。 & 背景损失 \\
\midrule
\multicolumn{4}{l}{\textbf{校准、反事实与场景控制}} \\
$consumer\_scale$ & $0.60$ & 固定缩放四大家鱼初值，保留既有物种占比，不由 2025 年观测标定估计。 & 初值规范化 \\
$gain\_scale$ & $1.45$ & 由 2025 年“四大家鱼恢复到禁渔前约 1.8 倍”这一总量约束单参数搜索得到。 & 唯一标定参数 \\
$F$（禁渔情景） & $0$ & 禁渔基线情景下共同捕捞压力取零；这里给出的是场景取值，而非新增独立参数。 & 基线政策设定 \\
$F_0$ & $0.25$ & 不禁渔反事实情景的固定共同捕捞压力，取自区间 $[0.10,0.40]$ 的中位值，不参与观测校准。 & 反事实对照 \\
$u$ & $0$ & 问题一中污染强度默认取零，避免与问题五的污染情景混用。 & 基线环境设定 \\
\midrule
\multicolumn{4}{l}{\textbf{观测代理与初值模板}} \\
$\mathrm{CPUE}^*$ 权重 & $(1,1,1,0.7)$ & 草鱼、鲢鱼、鳙鱼等权；青鱼因偏底栖、单网次可见性较低而保守下调。 & 捕获量代理 \\
$y_0^{\text{base}}$ & \makecell[l]{$(0.88,0.82,0.76,$\\$0.70,0.46,0.42,$\\$0.38,0.34)$} & 归一化初值模板，资源四维在前、鱼类四维在后；四类鱼在仿真前统一乘以 $consumer\_scale=0.60$。 & 状态初始化 \\
\bottomrule
\end{tabularx}
\normalsize
\end{table}
